\documentclass[11pt]{article}
\usepackage[T1]{fontenc}
\usepackage{lmodern,amsmath,amssymb,amsthm,mathtools}
\usepackage[margin=1in]{geometry}
\usepackage{microtype,booktabs,longtable,array,enumitem,xcolor}
\usepackage{hyperref}
\usepackage{xurl}
\hypersetup{colorlinks=true,linkcolor=blue!45!black,urlcolor=blue!50!black,
  pdftitle={Calibrated determinant guards for Diophantine ZFC consistency},
  pdfauthor={Prepared for Vladimir Reshetnikov; AI-assisted research}}
\newtheorem{theorem}{Theorem}[section]
\newtheorem{lemma}[theorem]{Lemma}
\newtheorem{proposition}[theorem]{Proposition}
\theoremstyle{definition}\newtheorem{definition}[theorem]{Definition}
\theoremstyle{remark}\newtheorem{remark}[theorem]{Remark}
\newcommand{\Z}{\mathbb Z}
\newcommand{\N}{\mathbb N}
\newcommand{\Pos}{\operatorname{Pos}}
\newcommand{\Con}{\operatorname{Con}}
\newcommand{\RefCode}{\operatorname{Ref}}
\newcommand{\ZFC}{\mathsf{ZFC}}
\newcommand{\sgn}{\operatorname{sgn}}
\newcommand{\PhiG}{\mathcal G}
\newcommand{\pin}{\texttt{9cca523d9f27ad39d52993f4065cbd7e58f8bdd4}}
\setlength{\emergencystretch}{2em}
\setlist{itemsep=3pt,topsep=5pt}
\title{Calibrated determinant guards\\for Diophantine ZFC consistency\\[5pt]
  \large A quartic in 38 integers, a sextic in 32, an octic in 27,\\
  and two positive-variable improvements}
\author{Prepared for Vladimir Reshetnikov and the ProveIt repository\\
  \small AI-assisted mathematical research; exact finite certificates included}
\date{23 October 2026}

\begin{document}
\maketitle
\begin{abstract}
The principal result is a new way to remove positivity witnesses without
raising the degree of a Diophantine equation. The canonical native encoding
already fixes the leading digits of two packed integers. Fixed rescalings
place those canonical witnesses in a narrow region where one quadratic
determinant is positive. Positivity of that determinant, for arbitrary
signed assignments, recovers both common-sign conditions previously imposed
separately. Replacing two ternary positivity representations by one reduces
the inspected 41-integer quartic to \textbf{38 unrestricted integer variables}
and 25 quadratic residuals. Exact coordinate eliminations give a sextic in
\textbf{32 integers} and an octic in \textbf{27 integers}.

Two independent companions give twelve positive variables at degree
\textbf{110} in \textbf{102} binary operations, and nine positive variables
at degree \textbf{1,379,730} in \textbf{478} operations. Every coefficient is
fixed by supplied finite data. The new proofs compose with explicitly named
native-compiler, Pell and syntactic-refutation results in ProveIt to give the
requested ZFC-provable unsolvability/consistency equivalences. No consistency
assumption, global optimum, literature-priority claim, or new end-to-end
proof-assistant certificate is asserted.
\end{abstract}
\clearpage
\tableofcontents
\clearpage

\section{What is constructed and what improves}
\label{sec:results}

The target is an ordinary polynomial with fixed integer coefficients,
not an exponential equation and not a formula with hidden quantifiers.
The convention for consistency is syntactic:
\[
 \Con(\ZFC)\;\equiv\;\neg\exists p\,\RefCode_{\ZFC}(p),
\]
where $\RefCode_{\ZFC}$ is the fixed recursive refutation predicate used by the
repository's calibrated endpoint. The desired theorem is an equivalence
\emph{provable in ZFC}; it does not ask ZFC to prove either consistency or
the absence of zeros separately.

The immutable comparison snapshot is \pin. In particular, the comparison
includes all eight wave-9 packages 60--67. The collection README's older
frontier is insufficient for this comparison. The main previous integer
points are package 64's 41-variable quartic and package 62's 35-variable
sextic and 30-variable octic \cite{p64,p62}. Table~\ref{tab:frontier}
records the changes, with different witness domains kept explicit.

\begin{table}[ht]
\centering\small
\begin{tabular}{@{}lrrrrl@{}}
\toprule
Domain & Unknowns & Degree & Operations & Prior point & Definition\\
\midrule
$\Z$ & \textbf{38} & 4 & -- & 41 / 4 & $Q_{38}$\\
$\Z$ & \textbf{32} & 6 & -- & 35 / 6 & $S_{32}$\\
$\Z$ & \textbf{27} & 8 & -- & 30 / 8 & $O_{27}$\\
$\Z_{>0}$ & 12 & \textbf{110} & \textbf{102} & 118 / 104 ops & $P_{12}$\\
$\Z_{>0}$ & 9 & \textbf{1,379,730} & \textbf{478} & 1,494,006 / 487 ops & $F_9$\\
$\N$ & 9 & 1,379,730 & 485 & separate translation & $F_{9,\N}$\\
\bottomrule
\end{tabular}
\caption{New explicit equations against the inspected wave-9 frontier.
The first three prior points are unknowns / degree. Operation counts for
the integer forms have not been optimized.}\label{tab:frontier}
\end{table}

The variable reductions are $3/41\approx7.32\%$, $3/35\approx8.57\%$, and
$3/30=10\%$, respectively. The nine-positive companion reduces the best
inspected degree by 114,276, or approximately $7.65\%$, and saves nine
operations relative to that same predecessor. Package 62's 453-operation
nine-positive equation still has a lower operation count at higher degree.
The ten-unrestricted-variable frontier is not improved.

There is no agreed scalar measure of ``simplest.'' In this report a degree
reduction, a variable reduction, and a shorter evaluator are separate
achievements. Binary arithmetic counts charge every $+,-,\times$, including
squarings and multiplication by fixed numerals; reuse and constant leaves
are free. They are not bit-complexity estimates. All quoted degrees count
only existential unknowns, after specializing the fixed coefficient ports.

\begin{theorem}[Constructed consistency equations]\label{thm:main}
With the complete coefficient bindings in the archive, the constructions
below satisfy, by composition with the named inherited compiler and
calibration theorems,
\begin{align*}
 \ZFC&\vdash (\neg\exists\mathbf x\in\Z^{38}\ Q_{38}(\mathbf x)=0)
                      \leftrightarrow\Con(\ZFC),\\
 \ZFC&\vdash (\neg\exists\mathbf x\in\Z^{32}\ S_{32}(\mathbf x)=0)
                      \leftrightarrow\Con(\ZFC),\\
 \ZFC&\vdash (\neg\exists\mathbf x\in\Z^{27}\ O_{27}(\mathbf x)=0)
                      \leftrightarrow\Con(\ZFC).
\end{align*}
The analogous assertions hold for $P_{12}$ over $\Z_{>0}^{12}$, $F_9$ over
$\Z_{>0}^9$, and its explicitly shifted companion over $\N^9$.
\end{theorem}

The proof dependency is substantive: this article gives complete ordinary
proofs of its new reductions, and imports the established native arithmetic
and refutation calibration used by the cited packages. It does not replace
those dependencies by numerical tests. Section~\ref{sec:scope} separates
that mathematical composition from the repository's still distinct task
of exporting a checked derivation in its object-language ZFC calculus.

\section{The new principle: one determinant for two sign conditions}
\label{sec:guard}

For integer variables define
\[
 \Pos(x,y,z)=1+x^2+y^2+z(z+1).
\]
It is always positive, since $z(z+1)\ge0$ for every integer $z$.
Conversely every positive integer $n$ is a value of $\Pos$: Legendre's
three-square theorem applies to $4n-3\equiv1\pmod4$, which is not of the
excluded form $4^a(8b+7)$. Exactly one of the three squares is odd. Writing
that term as $(2z+1)^2$ and the other two as $(2x)^2,(2y)^2$ gives the
claimed representation.

The previous quartic used two independent triples to impose
\begin{equation}\label{eq:oldguards}
 F\alpha=\Pos(f_1,f_2,f_3),\qquad
 FC=\Pos(c_1,c_2,c_3),\qquad q=F+\alpha+C.
\end{equation}
These rows require six private integer witnesses. The new idea is to use
the special location of \emph{canonical} accepting witnesses to choose a
stronger quadratic condition that requires only one triple.

\begin{lemma}[Weighted determinant guard]\label{lem:guard}
Let $A_0,B_0$ be fixed integers satisfying
$A_0,B_0\ge2$. For integers $q,F,C$ put
\[
 U=A_0F,\quad V=B_0C,\quad W=q-U-V,
 \qquad\PhiG=UV-W^2.
\]
If $\PhiG>0$, then $q,F,C,q-F-C$ are nonzero and have the same sign.
\end{lemma}
\begin{proof}
$UV>W^2\ge0$ implies that $U,V$ have one common nonzero sign.
Set $L=|U|+|V|$. The elementary inequality
$2\sqrt{|U||V|}\le L$ gives $|W|<L/2$. Thus $q=U+V+W$
has the same sign as $U,V$, and
\[
 |q|>L/2=(A_0|F|+B_0|C|)/2
                   \ge |F|+|C|.
\]
So $q,F,C$ have one sign and $|q|>|F|+|C|$, which gives the same sign
to $q-F-C$. The proof does not assume $q>0$.
\end{proof}

Consequently the single quadratic row
\begin{equation}\label{eq:newguard}
 \boxed{A_0B_0FC-(q-A_0F-B_0C)^2=\Pos(t_1,t_2,t_3)}
\end{equation}
implies both positivity conditions in~\eqref{eq:oldguards}. One may view
the left side as the determinant of
$\left(\begin{smallmatrix}U&W\\W&V\end{smallmatrix}\right)$.
The linear relation to $q$ recovers the common sign of this definite matrix.

Soundness alone is easy; completeness is the actual research issue.
Arbitrarily strengthening a positivity condition can destroy every witness.
Here the fixed native coefficient geometry supplies a uniform region
near $U=V=q/2$, in which the determinant remains positive.

\begin{lemma}[Uniform canonical region]\label{lem:region}
Let $s>0$, $q=2s$, and
$s\le U<8s/5$, $s\le V<6s/5$. Then
\[
 UV-(q-U-V)^2>\frac{12}{25}s^2>0.
\]
\end{lemma}
\begin{proof}
Write $U=s+x,V=s+y$, with $0\le x<3s/5$ and $0\le y<s/5$.
Then $W=-x-y$ and
\[
 UV-W^2=s^2+s(x+y)-x^2-xy-y^2
               >s^2-(9+3+1)s^2/25.
\]
\end{proof}

The reduction is therefore asymmetric in the useful way. An arbitrary
new zero restores an old zero by Lemma~\ref{lem:guard}; a refutation supplies
fresh canonical witnesses that satisfy Lemma~\ref{lem:region}. Equality of
the two zero sets on the same old coordinates is neither asserted nor
needed.

\section{Why the actual native encoding lies in that region}
\label{sec:canonical}

This section fixes the coefficients, rather than treating the guard as a
family with unbound parameters. Use the \emph{whole} source-34 tuple
\texttt{mixed-parity-unrounded-1064921} \cite{p34}. Write
\[
 \rho=2^{121},\quad L_0=8801,\quad B=\rho^{L_0}=2^{1064921},
 \qquad\beta=B-1.
\]
The first native tile, with ordered labels $(2,435,3,1)$, has clause
coefficient
\begin{equation}\label{eq:c0}
 c_0=105534389824824194647382298625,
 \qquad 2^{96}\le c_0<c_0+1\le3\cdot2^{95}<2^{97}<\rho.
\end{equation}
The two new fixed numerals are simply powers of two:
\begin{equation}\label{eq:scales}
 \boxed{A_0=2^{1113},\qquad B_0=2^{666104}.}
\end{equation}
In particular $A_0,B_0\ge2$. These are additional coefficients,
not existential variables. The archive reconstructs all six ports and
checks the four inherited raw values against their pinned SHA-256 hashes.

\subsection{Leading digits of the canonical \texorpdfstring{$C$}{C} word}

A finite accepting computation supplies a padded cyclic native word with
$N$ occupied cells and $q=B^N$. At each cell $j$, the source construction
puts one genuine tile, its occupied edge copies, and its four occupancy
anchors. The largest support exponent of its native integer $C_j$ is 3296,
with coefficient one. The next possible exponent is 2195. Therefore
\begin{equation}\label{eq:Cblock}
 \rho^{3296}\le C_j<\rho^{3296}+\rho^{2196}
                         \le(1+1/\rho)\rho^{3296}.
\end{equation}
The rank-tuning dummy is at exponent 1063; changing its permitted value
cannot affect either leading-support fact. The inequalities hold for every
size of a canonical accepting rectangle, including after all dummy changes.

\subsection{Leading digits of the canonical \texorpdfstring{$F$}{F} word}

Let $D_C,D_R$ be the fixed native convolution coefficients, so
$K_0=D_C+BD_R$ and $D_R=\rho$. The rank controller gives a whole-cell
temporal shift $w_c$ with
\[
 R\equiv1064921w_c\pmod{1064921N}.
\]
Thus the canonical transport word is the residue
\[
 F\equiv(D_C+BD_R+B^{w_c})C\pmod{q-1}.
\]
The actual block at cell $j$, with cyclic indices, is
\begin{equation}\label{eq:Fblockexact}
 F_j=D_CC_j+D_RC_{j-1}+C_{j-w_c}.
\end{equation}
The largest exponent in $D_C$ is 5495, with coefficient $c_0$.
The unique pair reaching exponent $8791=5495+3296$ in $D_CC_j$ has
coefficients $c_0$ and one. The other two summands in~\eqref{eq:Fblockexact}
have exponents at most 3297 and 3296.

It is necessary to exclude a carry into that leading coefficient. Let
$S=\sum_e c_e$ be the sum of the source clause coefficients. The sum of
coefficients of $D_C$ is $4+2S$. Each canonical $C_j$ coefficient is at
most three; the possible larger dummy value is still only three. Hence
each convolution coefficient in~\eqref{eq:Fblockexact} is at most
\[
 3(4+2S)+6=6S+18<991(2S+6)=\mathcal M<\rho.
\]
The last inequality is among the reconstructed source's strict finite
carry bounds; the stronger $4\mathcal M+8<\rho$ is checked. Thus no carry
changes the top digit, and
\begin{equation}\label{eq:Fblock}
 c_0\rho^{8791}\le F_j<(c_0+1)\rho^{8791}.
\end{equation}
Every block is strictly below $B-1$ since $8791<8801$ and $c_0<\rho$.
Their base-$B$ word is therefore the canonical representative in $(0,q-1)$;
the residue calculation has not confused zero with $q-1$.

\subsection{One estimate for every accepting witness length}

Put $J_N=(q-1)/(B-1)=\sum_{j<N}B^j$. The exponents were chosen because
\[
 121\cdot8791+1113=1064921-97,\qquad
 121\cdot3296+666104=1064921-1.
\]
Summing~\eqref{eq:Cblock} and~\eqref{eq:Fblock} therefore gives
\[
 \frac{c_0}{2^{97}}BJ_N\le A_0F<\frac{c_0+1}{2^{97}}BJ_N,
 \qquad
 \frac12BJ_N\le B_0C<\frac12(1+1/\rho)BJ_N.
\]
For all $N\ge1$,
$q\le BJ_N<qB/(B-1)$. Using~\eqref{eq:c0}, $\rho\ge8$ and $B\ge16$ gives
\[
 \frac q2\le A_0F<\frac34\frac{16}{15}q=\frac45q,
 \qquad
 \frac q2\le B_0C<\frac12\frac98\frac{16}{15}q=\frac35q.
\]
Set $s=q/2$. Lemma~\ref{lem:region} applies. This proves completeness
of the new determinant guard for the actual coefficient tuple. The large
sizes of $A_0,B_0$ express a known disparity of native digit positions; they
do not encode the answer to the consistency question.

\section{The complete 38-integer quartic}
\label{sec:quartic}

All variables in this section range over \emph{unrestricted integers}.
There are 32 arithmetic variables
\[
\begin{gathered}
 J,F,\alpha,z,f,h,r,v,w,\tau,\eta,\zeta,y,\gamma,o,j,\\
 N,E,P,c,a,X,D,C_2,L,F_2,V,g,I,B_a,\Delta,u
\end{gathered}
\]
and six guard variables $t_1,t_2,t_3,e_1,e_2,e_3$, totaling 38.
In the machine-readable packets, $v,w,B_a$ are named
\texttt{v\_scale}, \texttt{w\_scale}, \texttt{Baux}.
For readability define the polynomial abbreviations
\begin{align*}
 q&=\beta J+1,&C&=q-F-\alpha,&k&=\eta+\zeta,&H&=4a+3.
\end{align*}
\begin{align*}
 \PhiG&=A_0B_0FC-(q-A_0F-B_0C)^2,\\
 \mathcal R&=q(N-q-L)+F(q+L-N)+\alpha L+(M_C+qM_F)J.
\end{align*}
$q,C,k,H,\mathcal R,\PhiG$ are not additional unknowns.
Every displayed residual in the following list has degree at most two:
\begin{align*}
 r_1&=L-q^2+1,&r_2&=N-q(L+1),\\
 r_3&=X-4wq,&r_4&=a-E-(2v+1)N,\\
 r_5&=E(X+1)-Xa,&r_6&=c-\eta-(a-E)k,\\
 r_7&=P-E(c-\eta)-k,&r_8&=\tau^2-P^2+Pk-1,\\
 r_9&=D-X-ac-\gamma H,&r_{10}&=\Delta-a^2-4a-3,\\
 r_{11}&=C_2-c^2,&r_{12}&=D^2-\Delta C_2-1,\\
 r_{13}&=k-1-hE-\mathcal R,&
 r_{14}&=(K_0+4w)C-F-z(q-1),\\
 r_{15}&=F_2-f^2,&r_{16}&=V+c-fo,\\
 r_{17}&=V-jc+k-1-hE,&r_{18}&=V^2+g-y^2,\\
 r_{19}&=I-(2r+1)C_2,&r_{20}&=B_a-I^2,\\
 r_{21}&=B_a-\Delta(F_2-1),&r_{22}&=y^2-B_ag-1,\\
 r_{23}&=X-2u^2,&r_{24}&=\PhiG-\Pos(t_1,t_2,t_3),\\
 r_{25}&=\eta\zeta-\Pos(e_1,e_2,e_3).&&
\end{align*}
The promised single equation is
\begin{equation}\label{eq:Q38}
 \boxed{Q_{38}=\sum_{\nu=1}^{25}r_\nu^2=0.}
\end{equation}
The use of a sum is just a compact specification of one polynomial.
All 25 residuals and its expanded symbolic coefficient list are included.
The coefficient of $\tau^4$ is exactly one, so its degree is exactly four
on the actual fixed coefficient slice.

\subsection{Existence equivalence}

At an integer zero of $Q_{38}$ every residual is zero. The new guard gives
$F\alpha>0$ and $FC>0$ by Lemma~\ref{lem:guard}. Represent these two
positive integers by the old six private Pos variables. All other rows
are literally the corresponding package-64 quartic rows. Thus an integer
zero of $Q_{38}$ supplies a zero of that package's $Q_{41}$ and, by its
native/Pell soundness and K215 calibration, a finite ZFC refutation.

Conversely a refutation supplies the canonical witnesses in package 64's
completeness proof. They use the whole source-34 tuple, and satisfy all
retained rows, including $X=2u^2$ and the divisible-by-four scale.
Section~\ref{sec:canonical} proves $\PhiG>0$ for those very witnesses,
after rank tuning. A ternary Pos representation supplies the new three
variables. Hence $Q_{38}$ has an integer zero.

The quantifier order matters: one first obtains an accepting computation,
then a canonical native witness, then the new Pos representation. This
is why restricting the old numerical witnesses is legitimate. No search
for the absence of a witness, and no assumption of $\Con(\ZFC)$, occurs.

\subsection{Coefficient and proof costs}

The four old native numerals occupy 532,479 bytes in minimal unsigned
little-endian form. The two new leaves add 83,404 bytes, for a total of
615,883 bytes. The archive also provides short power-of-two recipes
and compressed payloads. This is a real coefficient-storage tradeoff for
three fewer variables at unchanged degree. It is not a claim that the
expanded polynomial's total bit length decreases. In fact the $J^2$
coefficient of the new residual is $-\beta^2(B_0-1)^2$, so squaring
that residual contributes $\beta^4(B_0-1)^4$ to the $J^4$ coefficient.
Including the retained rows, the exact coefficient is
\[
 [J^4]Q_{38}=\beta^4(B_0-1)^4+\beta^4+\beta^2(\beta-M_F)^2.
\]
It has exactly 6,924,100 bits for the supplied tuple, as checked both by
exact inequalities and by reconstruction. This is a lower bound on the
maximum coefficient bit length, not a claim that this coefficient is the
largest. Variable count, coefficient height and arithmetic description
length must therefore be assessed separately.

The inherited signed Pell argument remains unchanged. In particular the
new proof does not assume a positive radix before the existing argument
recovers its sign. It also preserves the corrected transport convention
$(K_0+4w)C=F+z(q-1)$; a different earlier convention uses $z+1$.

\section{Exact sextic and octic reductions}
\label{sec:elimination}

These reductions are purely algebraic. Whenever a residual is
$x-f(\text{other variables})$, deleting it and substituting $f$ everywhere
gives a bijection of integer zero sets, with the explicit polynomial
inverse restoring $x$. No divisibility or sign premise is needed.
The verifier checks the identity at each step and also substitutes the
complete final inverse into every original quartic row.

\subsection{A sextic in 32 integers}

First make the integral triangular change
\[
 \tau_{\rm old}=\tau_{\rm tail}+P-k.
\]
Then eliminate $P=k+E(c-\eta)$. The first Pell residual becomes
\[
 \tau_{\rm tail}^2+(2\tau_{\rm tail}-k)E(c-\eta)-1,
\]
which has degree three: its degree-four terms have cancelled exactly.
Eliminate, in addition,
\[
 F_2=f^2,\quad g=y^2-V^2,\quad X=2u^2,
 \quad N=q(L+1),\quad C_2=c^2.
\]
This removes six variables and six defining equations. The resulting 19
residuals have degree at most three in 32 integers. Let $S_{32}$ be their
sum of squares. The full inverse includes
\[
 \tau_{\rm old}=\tau_{\rm tail}+E(c-\eta).
\]
Only the retained strong-definition row
$I-(2r+1)c^2$ contains $r$. Its square gives
\[
 [r^2c^4]S_{32}=4,
\]
so $\deg S_{32}=6$ for the authentic fixed tuple. The full literal list,
expanded polynomial and inverse map are in
\path{integer_forms/sextic32.json} and its companion files.

\subsection{An octic in 27 integers, written out completely}

This version removes 11 registers. It is short enough to state without
any implicit source reference. Its 21 arithmetic unknowns are
\[
 J,F,\alpha,z,f,h,r,v,w,\tau,\eta,y,\gamma,o,j,E,c,a,V,I,u,
\]
and its six guard unknowns are again $t_1,t_2,t_3,e_1,e_2,e_3$.
Define only polynomial abbreviations
\begin{align*}
 q&=\beta J+1,&C&=q-F-\alpha,&k&=1-V+jc+hE,\\
 H&=4a+3,&\Delta&=a^2+H,&B_1&=E(c-\eta),\\
 R&=((q-1)(q-F)+\alpha)(q^2-1)+(M_C+qM_F)J,\\
 \PhiG&=A_0B_0FC-(q-A_0F-B_0C)^2.
\end{align*}
Here $B_1$ is a polynomial abbreviation and is unrelated to the fixed
native radix $B$. The 14 residuals are
\begin{align}
 s_1&=2u^2-4wq,\nonumber\\
 s_2&=a-E-(2v+1)q^3,\nonumber\\
 s_3&=E(2u^2+1)-2u^2a,\nonumber\\
 s_4&=c-\eta-(a-E)k,\nonumber\\
 s_5&=\tau^2-(k+B_1)B_1-1,\nonumber\\
 s_6&=(2u^2+ac+\gamma H)^2-\Delta c^2-1,\nonumber\\
 s_7&=-V+jc-R,\nonumber\\
 s_8&=(K_0+4w)C-F-z(q-1),\nonumber\\
 s_9&=V+c-fo,\nonumber\\
 s_{10}&=I-(2r+1)c^2,\nonumber\\
 s_{11}&=I^2-\Delta(f^2-1),\nonumber\\
 s_{12}&=y^2-I^2(y^2-V^2)-1,\nonumber\\
 s_{13}&=\PhiG-\Pos(t_1,t_2,t_3),\nonumber\\
 s_{14}&=\eta(k-\eta)-\Pos(e_1,e_2,e_3).\label{eq:O27rows}
\end{align}
The explicit equation is
\begin{equation}\label{eq:O27}
 \boxed{O_{27}=s_1^2+\cdots+s_{14}^2=0,\qquad\mathbf x\in\Z^{27}.}
\end{equation}
This is one convenient answer if a low degree and unrestricted integer
variables are preferred over the nine-positive high-degree construction.

Restore the eliminated quartic coordinates by
\begin{align*}
 L&=q^2-1,&N&=q^3,&X&=2u^2,\\
 C_2&=c^2,&F_2&=f^2,&g&=y^2-V^2,\\
 B_a&=I^2,&D&=2u^2+ac+\gamma H,&\Delta&=a^2+4a+3,\\
 \zeta&=k-\eta,&P&=k+E(c-\eta).&&
\end{align*}
These are integer polynomials. Substitution makes each original row
identically zero or one of the 14 listed residuals, with a possible harmless
sign change. Conversely projection recovers the octic coordinates.
Every residual has degree at most four, and
\[
 [a^4f^4]O_{27}=1.
\]
Only the strong row $I^2-\Delta(f^2-1)$ can supply that monomial.
Thus $\deg O_{27}=8$, not just an upper bound. Both eliminated forms inherit
the existence equivalence of $Q_{38}$ by exact inverse maps.

\input{positive12/section.tex}
\input{nine_positive/section.tex}

\section{Proof status, alternatives, and integration}
\label{sec:scope}

\subsection{The precise proof boundary}

The new determinant lemma, its canonical coefficient bounds, and the
integer elimination maps are elementary unbounded arguments. The two
positive companions also include complete ordinary arguments, including
the order in which denominator, sign, rank and divisibility facts are
recovered. All use finite integer arithmetic, elementary real inequalities
or standard number-theoretic lemmas that can be formalized in ZFC.

The inherited statements are explicitly the following: the finite native
encoding is sound for blank K215 accepting computations; every finite halt
has the canonical padded and rank-tuned native witnesses; the cited Pell
and ratio interfaces preserve existence; and the calibrated machine has a
finite blank halt exactly when the fixed ZFC calculus has a finite
refutation. The new code does not prove these facts by running a bounded
machine search.

The active repository distinguishes a closed generic nine-variable Lean
equation from the optimized native circuits, for which the authentic native
calibration in Lean and a derivation in the object-language ZFC calculus
remain open \cite{active}. Here the native bridge is imported from the cited
packages as an ordinary mathematical proof. This archive supplies
\emph{no newly compiled end-to-end Lean or Rocq
certificate}, and no exported derivation in the repository's object-level
ZFC syntax. Ordinary mathematical composition and an object-language
proof artifact are different deliverables. The statement in
Theorem~\ref{thm:main} has the former evidence, together with exact finite
certificates for its explicit data.

Negating the proved existence equivalence gives the requested statement.
It does not prove consistency. If ZFC is consistent, the new equations
have no zeros in their stated domains; that conditional assertion does
not entail a separate proof of their unsolvability inside ZFC.

\subsection{What the finite checks establish}

The root verification command rebuilds the coefficient tuple and polynomial
packets and checks the following independent evidence:
\begin{itemize}
\item All 430 literal machine transitions reconstruct the exact 943-tile
  list, and all four old coefficient bytes match their pinned hashes.
\item The new scales, leading support, unique top convolution pair, and
  strict carry and balancing hypotheses hold for those actual integers.
\item The quartic, sextic and octic have exactly the stated live unknowns,
  residual counts, exact degrees and nonzero leading coefficients.
\item Every monic elimination and every complete inverse restoration is
  an exact symbolic integer-polynomial identity.
\item A separately written checker, without importing the primary sparse
  engine, repeats the integer identities, coefficients and literal formulas.
\item The twelve-positive circuit has the exact degree-110 certificate and
  102 live operations. The nine-positive component checks all 11,616 source
  coefficient comparisons, the 264-term majorant certificate, degree-defect
  identities and an exact shared-power rewrite.
\end{itemize}
Finite sign boxes and modular evaluations are labeled as transcription
diagnostics. They are not substituted for universal proofs. The checks
also run with Python's \texttt{-O} flag; acceptance does not rely on
removable assertions.

\subsection{Why direct proof encoding was not promoted}

The request allowed a route independent of K215 and Jones. That direction
was investigated. Carl and Moroz give an explicit Diophantine representation
of predicate-calculus provability and specialize it to finitely axiomatized
G\"odel--Bernays set theory \cite{carlmoroz}. Their stated construction
uses 3,639,528 positive witnesses before a four-square translation. It is
a concrete alternative endpoint, but is not competitive with this
repository's present size bounds. Merely replacing the words ``Turing
machine'' by ``proof checker'' would not supply a smaller polynomial.

The determinant reduction itself does not use Jones's universal polynomial
or Sun's variable reduction; its low-degree forms inherit the native Pell
source. The nine-positive companion retains the existing Jones-style
ratio reduction. Sun's 2026 ten-integer result \cite{sun2026} remains
relevant to the separate minimum-integer-variable frontier. No new
eight-positive or nine-unrestricted-integer construction was obtained.

\subsection{A useful next mathematical target}

The structural lesson is specific: when canonical encoded quantities have
fixed leading digits, one should first look for a low-degree definite-form
region containing all those canonical values. Several independent
positivity witnesses may then be replaced by one. Arbitrary source zeros
need only map soundly; completeness can regenerate canonical witnesses.
The present proof works because two input words have fixed leading digit
positions and coefficients, unaffected by all padding and rank control.

Extending that mechanism to other guards requires a new uniform canonical
range theorem. The ratio $\eta/(\eta+\zeta)$ in the Pell part, for example,
has not been given an adequate fixed-coefficient balancing theorem here.
Dropping its remaining product guard would be unjustified. Similarly, a
different native tuple needs a fresh leading-support and carry audit before
using these two scales. This is a concrete research direction, not an
additional claimed result.

\subsection{Repository use and reproduction}

The archive is an independent research package. It need not replace the
active formally checked equation. Its root commands are
\begin{quote}\ttfamily
python3 verify\_all.py\\
python3 -O verify\_all.py\\
python3 build\_report.py
\end{quote}
The mathematical replay uses Python's standard library. A TeX Live
installation is needed only to recompile the article. Each positive
companion keeps its own complete coefficient binding: its ports must not
be mixed with the main integer tuple or with each other. The repository
snapshot, source identities, licenses, notices, deterministic polynomial
packets and verification receipts are included. The next numbered package
directory can be chosen at incorporation time.

\appendix
\section{Finite definition of the main native coefficients}
\label{app:coeff}

This appendix, together with the literal 943-tile JSON file, specifies
every coefficient of the integer equations by bounded finite arithmetic.
The original builder and independent raw hashes are supplied unchanged.
Let the ordered tiles be $(x_{it})_{0\le i<943,\ 0\le t<4}$, with Start
at index zero and Forbidden at index one. Put
\[
 e_t(j)=o_t+2j,\quad(o_0,o_1,o_2,o_3)=(1020,1000,980,960),
 \qquad W_j=2^{10+2j}.
\]
The support is
\[
 \mathcal E=\{0,\ldots,942\}\cup\{e_t(j):t<4,j<10\}
 \cup\{1063,1064,1095,1096,2195,3296\}.
\]
For each tile set
\[
 c_i=1+\sum_{t=0}^{3}\sum_{j=0}^{9}
          \operatorname{bit}_j(x_{it}+512)W_{10t+j}
          +\sum_{j=40}^{43}W_j.
\]
Set $c_{e_t(j)}=W_{10t+j}$, set the coefficients at
1095,1096,2195,3296 to $W_{40},W_{41},W_{42},W_{43}$ respectively,
and set $c_{1063}=c_{1064}=0$. Also put
\[
 \mu=1022+\sum_{j=0}^{43}W_j,\quad
 Z_0=(2^{119}-1)\sum_{e=8792}^{8798}\rho^e+(2^{27}-1)\rho^{8799}.
\]
Then
\begin{align*}
 D_C={}&\rho^{40}+\rho^{41}+\rho^{1100}+\rho^{2200}
       +\sum_{e\in\mathcal E}c_e(\rho^{4393-e}+\rho^{5495-e}),\\
 D_R={}&\rho,\\
 M_C={}&B-1-\sum_{\substack{e\in\mathcal E\\e\ne1}}\rho^e-2\rho^{1063},\\
 M_{F,n}={}&4+\mu(\rho^{4393}+\rho^{5495})
     +\sum_{j=0}^{9}(\rho^{1000+2j}+\rho^{1021+2j})
     +\rho^{2195}+\rho^{3296}+Z_0.
\end{align*}
The four inherited ports are exactly
\[
 \beta=B-1,\qquad K_0=D_C+BD_R,\qquad M_C,
 \qquad M_F=M_{F,n}+B-1.
\]
In particular $M_F$ is \emph{already shifted}. The two extra ports are
\eqref{eq:scales}. The literals contain no undecided search or parameter
that depends on whether ZFC is consistent. The compiler's finite
geometry, parity and rank-control proof is retained separately; a tuple
with the same rough numerical bounds would not be an adequate substitute.

\begin{thebibliography}{99}
\bibitem{p34} ProveIt, package 34, \emph{A mixed-radix, parity-separated,
unrounded coefficient tuple}, coefficient proof and literal native data,
at commit \pin.
\url{https://github.com/VladimirReshetnikov/ProveIt/blob/9cca523d9f27ad39d52993f4065cbd7e58f8bdd4/Computability/DiophantineConsistency/con-zfc-diophantine-forms/34-integer14/coefficients/COEFFICIENT_PROOF.md}.
\bibitem{p64} ProveIt, package 64, \emph{Structural reductions for explicit
Diophantine consistency equations}, especially \texttt{integer41/PROOF41.md}
and \texttt{PROOF44.md}, same commit.
\url{https://github.com/VladimirReshetnikov/ProveIt/tree/9cca523d9f27ad39d52993f4065cbd7e58f8bdd4/Computability/DiophantineConsistency/con-zfc-diophantine-forms/64-structural-reductions}.
\bibitem{p62} ProveIt, package 62, \emph{Sign recovery and centered norms
for explicit ZFC consistency equations}, especially the exact sextic/octic
eliminations and shared-power rewrite, same commit.
\url{https://github.com/VladimirReshetnikov/ProveIt/tree/9cca523d9f27ad39d52993f4065cbd7e58f8bdd4/Computability/DiophantineConsistency/con-zfc-diophantine-forms/62-sign-and-norm}.
\bibitem{p61} ProveIt, package 61, \emph{Congruence carriers}, positive
13/12-variable proofs and signed-index lemma, same commit.
\url{https://github.com/VladimirReshetnikov/ProveIt/tree/9cca523d9f27ad39d52993f4065cbd7e58f8bdd4/Computability/DiophantineConsistency/con-zfc-diophantine-forms/61-congruence-carriers}.
\bibitem{p65} ProveIt, package 65, \emph{Shared taps and integral carriers},
source tables and adaptive main-carrier proof, same commit.
\url{https://github.com/VladimirReshetnikov/ProveIt/tree/9cca523d9f27ad39d52993f4065cbd7e58f8bdd4/Computability/DiophantineConsistency/con-zfc-diophantine-forms/65-shared-taps}.
\bibitem{p63} ProveIt, package 63, \emph{Recovering signs and absorbing
Pell quotients}, deflated seventh-order majorant, same commit.
\url{https://github.com/VladimirReshetnikov/ProveIt/tree/9cca523d9f27ad39d52993f4065cbd7e58f8bdd4/Computability/DiophantineConsistency/con-zfc-diophantine-forms/63-sign-recovery}.
\bibitem{p67} ProveIt, package 67, \emph{Oriented guards}, wave-9 degree
and operation comparison, same commit.
\url{https://github.com/VladimirReshetnikov/ProveIt/tree/9cca523d9f27ad39d52993f4065cbd7e58f8bdd4/Computability/DiophantineConsistency/con-zfc-diophantine-forms/67-oriented-guards}.
\bibitem{active} ProveIt, \emph{A Diophantine equation for syntactic
consistency of ZFC}, active formalization status, same commit.
\url{https://github.com/VladimirReshetnikov/ProveIt/blob/9cca523d9f27ad39d52993f4065cbd7e58f8bdd4/Computability/DiophantineConsistency/README.md}.
\bibitem{carlmoroz} M. Carl and B. Z. Moroz, \emph{On a Diophantine
representation of the predicate of provability}, MPIM preprint 2007-132;
published in Zap. Nauchn. Sem. POMI 407 (2012), 77--104, English translation
J. Math. Sci. 199 (2014), 36--52.
\url{https://webdoc.sub.gwdg.de/ebook/serien/e/mpi_mathematik/2007/132.pdf}.
\bibitem{sun2026} Zhi-Wei Sun, \emph{Ten unknowns for Hilbert's tenth
problem over the integers}, arXiv:2609.01594v2, September 2026.
\url{https://arxiv.org/abs/2609.01594}.
\end{thebibliography}
\end{document}

